\section{Conclusion}
\label{sec:conclusion}

This paper evaluated whether an explicit lifecycle for stored strategy revisions
improved an LLM inventory agent's performance after operating conditions
changed.
ShiftMem made reuse decisions explicit through conditional retrieval, delayed
validation, and an auditable lifecycle, while a shared deterministic controller
retained responsibility for daily orders. In the pre-Test-frozen evaluation, the
dependence-limited prespecified decision rule did not support a ShiftMem
advantage over the implemented lexical retrieval baseline. The mean point
estimate was unfavorable, and the post-hoc clustered sensitivity retained the
same direction.

The result does not establish equivalence or show that the conditional memory
lifecycle is intrinsically ineffective. Point estimates varied descriptively
across the two
model deployments and test regimes, and provider and structured-output failures
were retained as part of the observed system behavior. Neither the descriptive
heterogeneity nor the reliability associations provided mechanism-level
evidence: the associations were observational, key lifecycle ablations were not
run, and every pair executed the lexical retrieval baseline first. The evidence
is therefore limited to the tested synthetic single-item environment, provider
deployments, lexical retrieval baseline, and fixed execution order.

The main methodological outcome is a documented implementation and evaluation
architecture for operational memory under delayed feedback: bounded strategy
review, deterministic execution, retained failures, delayed validation,
content-addressed evidence, and network-free analysis verification. A stronger
follow-up should randomize or counterbalance method order, add independent
LLM-output replications, encode
demand as an applicability condition, and align the available history with the
forecast-window range. It should also compare an embedding-based retrieval
baseline and run the dormancy/reactivation and detector ablations required to
identify component-level effects. These controls are prerequisites for a broader
claim about conditional memory efficacy.
